% Task 1 (`Knowledge of the sky') of the three-task planetarium round.
The projector will display the sky as seen from near the equator
($0^\circ$\,N, $19^\circ$\,E). The rotation of the sky will be stopped for
about 2 minutes for part (a), then it will start to rotate for parts (b) and
(c). The objects for parts (b) and (c) will be displayed simultaneously.
(Projection time 10 minutes.)

\begin{description}
\item[(a)] A meteor shower will be visible in the sky. Determine the
  constellation of the radiant and estimate its right ascension and
  declination coordinates. \hfill(3 points)
\item[(b)] Identify which of the following variable stars visible in the sky
  are in low (write `DIM') or high (write `BRIGHT') brightness states. The
  mean magnitude as shown in the atlas and the magnitude range are given for
  each star.

  \begin{center}
  \begin{tabular}{lccc}
  Name & atlas mag. & mag. range & DIM / BRIGHT \\
  \hline
  $\gamma$ Cas (Cih) & 2 & 1.6 -- 3.0 & \\
  $\delta$ Cep & 4 & 3.5 -- 4.4 & \\
  $\mu$ Cep (Erakis) & 4 & 3.4 -- 5.1 & \\
  $\beta$ Per (Algol) & 2 & 2.2 -- 3.4 & \\
  $o$ Cet (Mira) & 3.5 & 2.0 -- 10.1 & \\
  $\chi$ Cyg & 4.5 & 3.3 -- 14.1 & \\
  L$_2$ Pup & 4.5 & 2.6 -- 6 & \\
  $\delta$ Sco (Dschubba) & 2 & 1.6 -- 2.3 & \\
  \end{tabular}
  \end{center}
  \hfill(8 points)
\item[(c)] Identify the constellations whose borders are marked and give
  their IAU abbreviations. \hfill(9 points)
\end{description}
